% ==== note-gen preamble (English notes). For Chinese, swap to xelatex + ctex. ====
\documentclass[11pt,a4paper]{article}
\usepackage[a4paper,top=2.5cm,bottom=2.5cm,left=2.5cm,right=2.5cm]{geometry}
\usepackage{amsmath,amssymb,amsthm}
\usepackage{graphicx}
\usepackage{float}
\usepackage{booktabs}
\usepackage{array}
\usepackage{xcolor}
\usepackage{tcolorbox}
\usepackage[colorlinks=true,allcolors=blue]{hyperref}
\usepackage{enumitem}
\usepackage{setspace}
\usepackage{microtype}
\usepackage{cite}

\onehalfspacing
\newtheorem{definition}{Definition}
\newtheorem{intuition}{Intuition}

% Highlighted "key idea" callout box
\newtcolorbox{keybox}[1][]{colback=blue!5!white,colframe=blue!60!black,
  fonttitle=\bfseries,title=Key Idea,#1}
\newtcolorbox{pitfall}[1][]{colback=red!5!white,colframe=red!60!black,
  fonttitle=\bfseries,title=Common Pitfall,#1}

\newcommand{\eg}{\textit{e.g.}}
\newcommand{\ie}{\textit{i.e.}}
